\documentclass[11pt]{article}
\usepackage{amsmath,amssymb,amsthm,mathtools,geometry,hyperref}
\geometry{margin=1in}
\newtheorem{theorem}{Theorem}
\newtheorem{proposition}{Proposition}
\newcommand{\C}{\mathbb C}
\newcommand{\Disc}{\operatorname{Disc}}
\newcommand{\Spec}{\operatorname{Spec}}
\newcommand{\Res}{\operatorname{Res}}
\newcommand{\cP}{\mathcal P}
\title{Complete fibres, the eight-sheet domain, and the nonproper-value locus of the ES--Fable map}
\author{Independent programme calculation}
\date{20 September 2026}
\begin{document}
\maketitle

The input polynomial is the four-dimensional map in the canonical
ES--Fable reader \cite[Theorem
\texttt{explicit-four-time-Jacobian-collision}]{ES488}.
Its complete inverse calculation below includes every multiple-root
and missing-chart case. The original conductor is retained as a
fixed linear receiving isomorphism in the final section; its
coefficients and its original Gamma norm are not replaced.

\section{The original polynomial and the factor coordinates}
Fix $i^2=-1$. The original coordinates are $(a,y,z,w)\in\C^4$,
and the ordered output coordinates are $(u_0,u_2,u_3,u_4)$. Define
\begin{align}
 P_0={}&a^3z+2a^2y-ia,\nonumber\\
 P_2={}&-a^3y^2z-2ia^2yz+a^2w-2a^2y^3-10iay^2+3az+y,\nonumber\\
 P_3={}&2a^3y^3z+6ia^2y^2z+2a^2wy+4a^2y^4+2iaw
       -4iay^3-2ayz+2iz+7y^2,\nonumber\\
 P_4={}&2a^3y^4z+8ia^2y^3z+a^2wy^2+4a^2y^5+2iawy
       +7iay^4-10ay^2z-4iyz-w-3y^3.
 \tag{GF1}
\end{align}
The factor-coordinate map $\Phi(a,y,z,w)=(a,b,c,d,e,f)$ is
\begin{align}
 b&=i+ay,&c&=-i+2ay+a^2z,\nonumber\\
 d&=-iy-a(iz+2y^2)-a^2yz,&e&=2z-7iy^2+aw,\nonumber\\
 f&=iw+3iy^3-4yz+a(6iy^2z+wy+4y^4)+2a^2y^3z.
 &&\tag{GF2}
\end{align}
Let
\[
 L=aU+bV,\qquad C=cU^3+dU^2V+eUV^2+fV^3.
\]
Their original incidence and resultant equations are
\begin{equation}
 ad+bc=1,\qquad
 R(L,C)=C(-b,a)=a^3f-a^2be+ab^2d-b^3c=1.\tag{GF3}
\end{equation}
Direct substitution of GF2 proves both identities. Multiplication
of the two factors gives the complete coefficient map
\begin{equation}
 LC=(ac)U^4+U^3V+(ae+bd)U^2V^2+(af+be)UV^3+(bf)V^4,
 \qquad P=(ac,ae+bd,af+be,bf).\tag{GF4}
\end{equation}
Substitution of GF2 in GF4 gives precisely GF1.

On $a\ne0$, every factor state satisfying GF3 has the unique inverse
\begin{equation}
 y=(b-i)/a,\qquad z=(c+i-2ay)/a^2,\qquad
 w=(e-2z+7iy^2)/a.\tag{GF5}
\end{equation}
These formulas recover $a,b,c,e$. The first equation GF3 uniquely
determines $d$, and its second uniquely determines $f$ because
$a^3\ne0$. Thus GF5 is a full inverse on that open set.
At $a=0$, GF3 forces $bc=1$ and $-b^3c=1$, hence $b=\pm i$.
The original chart GF2 has $b=i$, $c=-i$, and contains all such
states: its inverse there is
\[
 y=id,\qquad z=(e+7iy^2)/2,\qquad
 w=(f-3iy^3+4yz)/i.
\]
It contains none of the $a=0,b=-i$ states. This description of the
chart is used explicitly in every fibre count below.

\section{Exact inverse over every target}
For $u=(u_0,u_2,u_3,u_4)$ define
\begin{equation}
 H_u(U,V)=u_0U^4+U^3V+u_2U^2V^2+u_3UV^3+u_4V^4,
 \qquad h_u(x)=H_u(x,1).\tag{GF6}
\end{equation}
The fixed $U^3V$ coefficient is part of the original map, not a
choice made during this calculation.

\begin{theorem}[Complete fibre, including every multiple-root case]
Every finite simple root $r$ of $h_u$ contributes exactly two
distinct points of $P^{-1}(u)$. Their complete coordinates are
\begin{align}
 a^2&=1/h_u'(r),\qquad a\in\{\text{both square roots}},\nonumber\\
 y&=-r-i/a,\nonumber\\
 z&=u_0/a^3+2r/a+3i/a^2,\nonumber\\
 w&=7ir^2/a+(u_2-17r+u_0r^2)/a^2-13i/a^3-2u_0/a^4.
 \tag{GF7}
\end{align}
A finite multiple root contributes no point. When $u_0=0$ there
is exactly one additional point, namely
\begin{equation}
 q_\infty(u)=
 \left(0,u_2,\frac{i(7u_2^2-u_3)}2,
                  11u_2^3-2u_2u_3-u_4\right).\tag{GF8}
\end{equation}
There are no other points. In particular, writing $m(u)$ for the
number of distinct finite simple roots,
\begin{equation}
 \#P^{-1}(u)=2m(u)+\mathbf1_{\{u_0=0\}}.\tag{GF9}
\end{equation}
\end{theorem}
\begin{proof}
A source point with $a\ne0$ gives the marked finite root
$r=-b/a$ of $H_u=LC$. The quotient factor is necessarily
\begin{align*}
 b&=-ar,&c&=u_0/a,&d&=(1+u_0r)/a,\\
 e&=(u_2+r+u_0r^2)/a,&
 f&=(u_3+u_2r+r^2+u_0r^3)/a.
\end{align*}
These formulas follow by comparing successive coefficients in
$H_u=a(U-rV)C$. The final constant coefficient is equivalent to
$h_u(r)=0$. Evaluating $C$ at $(-b,a)=(ar,a)$ gives
\[
 R(L,C)=a^3 C(r,1)=a^2 h_u'(r).
\]
Thus GF3 holds exactly when $a^2h_u'(r)=1$. A multiple root is
excluded by this equation, while each simple root supplies exactly
two nonzero values of $a$. The other incidence equation follows
without a further constraint:
$ad+bc=1+u_0r-u_0r=1$. The complete inverse GF5 now gives GF7;
for example substitution in $w=(e-2z+7iy^2)/a$ yields all four
displayed terms of its last line. These operations are reversible,
so they prove both sufficiency and exhaustiveness on $a\ne0$.
Distinct roots give different $r=-b/a$, and the two signs of $a$
give distinct source points.

If $a=0$ then $u_0=ac=0$. Conversely, GF1 on this hyperplane reads
\[
 P(0,y,z,w)=(0,y,2iz+7y^2,-4iyz-w-3y^3).
\]
Solving successively gives exactly GF8. In projective factor terms,
$u_0=0$ makes $[U:V]=[1:0]$ a simple root: $H_u=V(U^3+\cdots)$
and the leading coefficient of the parenthesis is exactly one.
Its two resultant-one factor scalings have $b=\pm i$, but GF2
retains only $b=i$. This is precisely the single point just solved,
regardless of any multiple finite roots. Hence there are no omitted
chart exceptions, and GF9 follows.
\end{proof}

For clarity the exhaustive multiplicity counts are
\[
\begin{array}{c|c|c}
\text{target chart}&\text{finite root multiplicities}&\#P^{-1}(u)\\\hline
u_0\ne0&1+1+1+1&8\\
u_0\ne0&2+1+1&4\\
u_0\ne0&3+1&2\\
u_0\ne0&2+2\text{ or }4&0\\\hline
u_0=0&1+1+1&7\\
u_0=0&2+1&3\\
u_0=0&3&1
\end{array}
\]
Every row follows directly from GF9. The polynomial always has
degree four or three, because the cubic coefficient is fixed at one.

\section{A short complete Jacobian proof}
On $a\ne0$, use parameters $(a,r,A,B)$, where $A=u_0$ and $B=u_2$.
The two equations $h_u(r)=0$, $a^2h_u'(r)=1$ solve as
\begin{equation}
 u_3=a^{-2}-4Ar^3-3r^2-2Br,
 \qquad u_4=3Ar^4+2r^3+Br^2-r/a^2.\tag{GF10}
\end{equation}
GF7 gives the source parameterization
\[
 q(a,r,A,B)=\left(a,-r-i/a,A/a^3+2r/a+3i/a^2,
 7ir^2/a+(B-17r+Ar^2)/a^2-13i/a^3-2A/a^4\right).
\]
Its Jacobian determinant in this parameter order is $-a^{-5}$:
after its first row, the remaining triangular minor has diagonal
$-1,a^{-3},a^{-2}$. For the target parameterization
$(A,B,u_3,u_4)$, reorder the first two parameter columns past the
last two; this permutation has positive sign. The remaining
two-by-two determinant is
\[
 \det\begin{pmatrix}
 -2/a^3&-(12Ar^2+6r+2B)\\
 2r/a^3&r(12Ar^2+6r+2B)-a^{-2}
 \end{pmatrix}=2a^{-5}.
\]
The chain rule therefore proves
\begin{equation}
 \det DP=-2\quad\text{on }a\ne0.\tag{GF11}
\end{equation}
Every original source point with $a\ne0$ has these parameters by
GF5--7, so this identity holds on that entire open set. Both sides
of GF11 are polynomials in the original coordinates; their
difference vanishes on a nonempty dense open set and is consequently
the zero polynomial. This proves GF11 everywhere, including the
$a=0$ chart, with its original constant and sign.

\section{The exact omitted image}
The only quartic multiplicity patterns with no simple root are
$2+2$ and $4$. Equivalently a quartic with $u_0\ne0$ and no simple
root is the square of a quadratic times its leading coefficient.
Because its cubic coefficient is one, it must have the form
\[
 H_u(U,V)=u_0\left(U^2+\frac{1}{2u_0}UV+u_3V^2\right)^2.
\]
Indeed, write the monic quadratic as $U^2+pUV+qV^2$; comparison
of cubic and linear coefficients gives $p=1/(2u_0)$ and $q=u_3$.
The remaining coefficients then give the exact image complement
\begin{equation}
 \C^4\setminus P(\C^4)
 =\{4u_0u_2-1-8u_0^2u_3=0,\quad u_4-u_0u_3^2=0\}.
 \tag{GF12}
\end{equation}
The first equation excludes $u_0=0$. Conversely the two equations
give exactly the displayed square. GF9 therefore proves both
inclusions in GF12, not only exclusion of the square targets.

\section{Discriminant and the genuine eight-sheet domain}
Write $A=u_0,B=u_2,C=u_3,D=u_4$. Define the original binary-quartic
discriminant polynomial by
\begin{align}
 \mathfrak D(u)={}&256A^3D^3-128A^2B^2D^2+144A^2BC^2D
 -27A^2C^4-192A^2CD^2\nonumber\\
 &+16AB^4D-4AB^3C^2-80AB^2CD+18ABC^3
 +144ABD^2-6AC^2D\nonumber\\
 &-4B^3D+B^2C^2+18BCD-4C^3-27D^2.
 \tag{GF13}
\end{align}
For $A\ne0$, the Sylvester determinant of $h_u$ and $h_u'$,
divided by $A$, expands to GF13. If $r_1,\ldots,r_4$ are its
roots counted with multiplicity, that same determinant gives
\begin{equation}
 \mathfrak D=A^6\prod_{j<k}(r_j-r_k)^2.\tag{GF14}
\end{equation}
One can verify this identity first for distinct roots: the
resultant is $A^3\prod_jh_u'(r_j)$, and
$\prod_jh_u'(r_j)=A^4\prod_{j<k}(r_j-r_k)^2$ because
$(-1)^{4\cdot3/2}=1$. Polynomial continuity gives the identity
also at multiple roots. At $A=0$, GF13 becomes
\[
 B^2C^2-4C^3-4B^3D+18BCD-27D^2,
\]
the discriminant of the monic cubic $x^3+Bx^2+Cx+D$. The same
resultant calculation in degree three proves that it vanishes
exactly at a repeated finite root. Since the projective root at
infinity is simple when $A=0$, GF13 vanishes exactly when $H_u$
has any multiple projective root, on both charts.

Let
\begin{equation}
 \mathcal B=\{u\in\C^4:u_0\mathfrak D(u)\ne0\},\qquad
 R=\C[u_0,u_2,u_3,u_4,(u_0\mathfrak D)^{-1}].\tag{GF15}
\end{equation}
Over this precise domain the map $P$ is an eight-sheet finite
unramified cover. Its complete coordinate algebra is
\begin{equation}
 S=R[r,a]/(h_u(r),\ a^2h_u'(r)-1).\tag{GF16}
\end{equation}
Here is a direct proof of each assertion. Since $u_0$ is a unit,
the first relation is a monic relation of degree four after division
by $u_0$, and $R[r]/(h_u)$ is free of rank four over $R$ with basis
$1,r,r^2,r^3$. The resultant of $h_u$ and $h_u'$ is a unit in $R$.
The adjugate of their Sylvester matrix therefore gives a Bezout
identity over $R$, so $h_u'(r)$ is a unit in this rank-four algebra.
The second relation is consequently the monic relation
$a^2=h_u'(r)^{-1}$. Thus $S$ is free of rank eight over $R$, with
basis $1,r,r^2,r^3,a,ar,ar^2,ar^3$.

The original chart inverse GF7 has entries in $S$: $a$ is a unit
with inverse $a h_u'(r)$. In the other direction, for an original
source point over $\mathcal B$ the identity $u_0=ac$ makes $a$
a unit, and $r=-b/a$ is regular. GF3--7 prove these two maps are
mutually inverse. Thus GF16 is the coordinate algebra of the
\emph{actual} preimage $P^{-1}(\mathcal B)$, rather than a new
eight-point model.

The derivative matrix of the two defining equations of GF16 with
respect to $(r,a)$ has determinant
$2a\,h_u'(r)^2$, a unit. This proves that each of its eight simple
solutions gives one local holomorphic inverse, equivalently that
the finite free map is unramified. Explicitly the four simple roots
vary holomorphically in a sufficiently small target neighbourhood,
and each nonzero $1/h_u'(r)$ has its two separate holomorphic
square roots there. GF7 gives eight disjoint local inverse charts.
There are exactly eight by GF9. These statements require neither
a chosen global square-root branch nor identification of its two signs.

\section{The complete nonproper-value locus}
A target $u$ is called a nonproper value here if there is a sequence
$q_n\in\C^4$ escaping every compact set with $P(q_n)\to u$.
For this polynomial map that is the local failure of properness
near $u$. The exact locus is
\begin{equation}
 \boxed{\quad\mathcal N(P)=\{u_0=0\}\ \cup\ \{\mathfrak D(u)=0\}.\quad}
 \tag{GF17}
\end{equation}

First take $u\in\mathcal B$ and a convergent sequence of targets
$u_n\to u$. On a sufficiently small closed neighbourhood of $u$,
$u_0$ and $\mathfrak D$ are bounded away from zero. Every finite
root is bounded: after dividing $h$ by $u_0$, any root has modulus
at most one plus the maximum modulus of its nonleading
coefficients. This follows by comparing the leading term with
the finite geometric sum of the remaining terms outside that disk.
Moreover $|h_{u_n}'(r_n)|$ is bounded away from zero uniformly over
all roots. Otherwise a convergent subsequence of the bounded
$r_n$ would give a root $r$ of $h_u$ with $h_u'(r)=0$,
contradicting $\mathfrak D(u)\ne0$. The derivatives also have a
uniform upper bound, by the root and coefficient bounds. Hence
$a^2=1/h'(r)$ bounds both $a$ and $a^{-1}$; GF7 then bounds every
original source coordinate. There is no $a=0$ solution on this
neighbourhood. Thus no sequence above $u_n$ can escape, proving
$\mathcal N(P)\subseteq\C^4\setminus\mathcal B$.

For the reverse inclusion at $u_0=0$, let GF8 be
$q_0=(0,y_0,z_0,w_0)$ and set
$q_\epsilon=(\epsilon,y_0,z_0,w_0)$ for real $\epsilon>0$.
On $a\ne0$ the exact factor-sign involution is
\begin{equation}
 \Sigma(a,y,z,w)=\left(-a,y+\frac{2i}{a},\frac{6i}{a^2}-z,
 w-\frac{14iy^2}{a}+\frac{28y}{a^2}+\frac{40i}{a^3}\right).
 \tag{GF18}
\end{equation}
To verify its invariance, substitute GF18 in GF2: $a,b,c,e$
all change sign. GF3 then forces $d$ and $f$ also to change sign
because $a\ne0$. Both factors in GF4 change sign, so
$P\Sigma=P$. The full inverse GF5 proves $\Sigma^2=I$ as well.
Now the second coordinate of $\Sigma q_\epsilon$ is
$y_0+2i/\epsilon$ and escapes, while
$P\Sigma q_\epsilon=Pq_\epsilon\to Pq_0=u$.
Every point of the entire hyperplane $u_0=0$ is consequently
nonproper, including its discriminant intersections.

Finally let $u_0\ne0$ and $\mathfrak D(u)=0$. Choose any multiple
root $r_0$ of $h_u$ and define
\[
 h_\epsilon(x)=h_u(x)+\epsilon(x-r_0),\qquad
 u_\epsilon=(u_0,u_2,u_3+\epsilon,u_4-\epsilon r_0).
\]
These preserve the original fixed cubic coefficient.
The root $r_0$ persists and has derivative exactly $\epsilon$.
For real $\epsilon>0$ choose $a=\epsilon^{-1/2}$ and apply the
complete inverse GF7 to this simple root. It gives a genuine
source point $q_\epsilon$ with $Pq_\epsilon=u_\epsilon\to u$,
while its first coordinate tends to infinity. This argument works
for every multiplicity of $r_0$ and does not require the other roots
to be simple. It proves the remaining inclusion in GF17.

\section{The literal arithmetic quartet as a marked target quartic}
Retain the original four distinct roots, with both original signs:
\[
 \rho_{\varepsilon,\eta}=\tfrac12+\varepsilon\delta+i\eta\gamma,
 \qquad\varepsilon,\eta\in\{1,-1\},\quad
 0<\delta<\tfrac12,\quad\gamma>2.
\]
Their sum is exactly two. Therefore the quartic with these literal
roots and the required original cubic coefficient one is
\begin{equation}
 H_\rho(U,V)=-\tfrac12\prod_{\varepsilon,\eta}
                    (U-\rho_{\varepsilon,\eta}V).\tag{GF19}
\end{equation}
The factor $-1/2$ is forced by the cubic coefficient and is retained
in every derivative, discriminant, and inverse below. With
$x=S-1/2$, its affine polynomial is
\[
 h_\rho(S)=-\tfrac12\bigl(x^4+2(\gamma^2-\delta^2)x^2
                              +(\delta^2+\gamma^2)^2\bigr).
\]
Expanding gives the exact target
\begin{equation}
 u_\rho(\delta)=\left(-\tfrac12,\
 \delta^2-\gamma^2-\tfrac34,\
 \gamma^2-\delta^2+\tfrac14,\
 -\tfrac1{32}-\frac{\gamma^2-\delta^2}{4}
                     -\frac{(\delta^2+\gamma^2)^2}{2}\right).
 \tag{GF20}
\end{equation}
At the marked root $\rho_{\varepsilon,\eta}$,
\begin{equation}
 h_\rho'(\rho_{\varepsilon,\eta})
       =4\delta\gamma(\varepsilon\gamma-i\eta\delta),\qquad
 a_{\varepsilon,\eta,\pm}^{2}
       =\frac1{4\delta\gamma(\varepsilon\gamma-i\eta\delta)}.
 \tag{GF21}
\end{equation}
To verify the derivative, the three differences from that root are
$2i\eta\gamma$, $2\varepsilon\delta$, and
$2\varepsilon\delta+2i\eta\gamma$; multiply them and the retained
leading coefficient $-1/2$. This gives precisely GF21. Both square
roots are present at every one of the four marked roots.

In the order $(+,+),(+,-),(-,+),(-,-)$ the original Vandermonde is
$\Delta_\rho=-64\delta^2\gamma^2(\delta^2+\gamma^2)$. Hence
\begin{equation}
 \mathfrak D(u_\rho)=\Delta_\rho^2/64
                   =64\delta^4\gamma^4(\delta^2+\gamma^2)^2>0.
 \tag{GF22}
\end{equation}
The denominator $64$ is the sixth power of the retained factor
$2$. Thus $u_\rho\in\mathcal B$, and GF7 and GF21 give all
eight original chart states over it. Unlike the earlier target
$UV(U-V)(U+V)$, it has no root at infinity and loses no sign
partner from the source chart.

For fixed $\gamma$ and $\delta\downarrow0$, GF19 tends to
\begin{equation}
 H_0=-\tfrac12\left((U-V/2)^2+\gamma^2V^2\right)^2.\tag{GF23}
\end{equation}
This target is on the omitted square surface GF12, so its complete
source fibre is empty. The eight nearby states all escape; they
do not become fewer finite chart points. Indeed GF21 gives
\[
 |a|=\bigl(4\delta\gamma\sqrt{\gamma^2+\delta^2}\bigr)^{-1/2},
 \qquad \sqrt\delta\,|a|\longrightarrow\frac1{2\gamma}.
\]
For each continuous choice of its sign branch,
\begin{equation}
 \sqrt\delta\,a_{\varepsilon,\eta,\pm}\longrightarrow
 \kappa_{\varepsilon,\pm},\qquad
 \kappa_{\varepsilon,\pm}^2=\frac1{4\varepsilon\gamma^2}.
 \tag{GF24}
\end{equation}
GF7 then gives
$y\to-(1/2+i\eta\gamma)$ and $z,w\to0$. Consequently, in
the original source coordinates,
\begin{equation}
 \sqrt\delta\,q_{\varepsilon,\eta,\pm}(\delta)
       \longrightarrow\kappa_{\varepsilon,\pm}e_a,
 \qquad e_a=(1,0,0,0)^{\mathsf T}.\tag{GF25}
\end{equation}
Each claim follows from the displayed inverse formulas; no moving
period or limiting identification of distinct root coordinates is
used to infer these limits.

\section{All eight escaping branches in one fixed original conductor norm}
Fix once and for all an admitted reference quartet with parameter
$\delta_*>0$, its original conductor, an original permitted source
order $s$, and its original target Gamma norm on $W=\cP_3$.
Let $\Psi_*:\C^4\to W$ be the exact invertible receiving map
constructed from that reference conductor in FC26--35
\cite{Bridge}. Its entries include the reference period and every
reference moment; none is varied in this section. Define
\[
 K=\begin{pmatrix}
 0&i&1/\sqrt2&1/\sqrt2\\
 0&-1&-1-\sqrt2 i&1-\sqrt2 i\\
 i/2&-3i&2\sqrt2+6i&-2\sqrt2+6i\\
 0&13&-34-19\sqrt2 i&34-19\sqrt2 i
 \end{pmatrix},\quad \det K=77\sqrt2 i/2.
\]
For complete coordinates, let $V_*=(\rho_{*,j}^n)_{0\le n\le3,1\le j\le4}$
be the reference root Vandermonde, and let $\mu_j$ and $v$ be
the original reference conductor moments and their actual first
nonzero order. In the target basis $(1,S',S'^2,S'^3)$ the complete
matrix of this fixed map is
\[
 [\Psi_*]=B_{v,*}(V_*^{\mathsf T})^{-1}K^{-1},\qquad
 (B_{v,*})_{kj}=\begin{cases}
 \binom{v+j}{k}\mu_{v+j-k},&0\le k\le j\le3,\\
 0,&k>j.
 \end{cases}
\]
The first determinant is
$\det B_{v,*}=\mu_v^4\prod_{j=0}^3\binom{v+j}{v}\ne0$;
the reference Vandermonde is nonzero and $K$ has the determinant
displayed above. This proves the asserted invertibility with all
original factors retained. The target norm in this monomial frame
is the original one, not the Euclidean norm of the column:
\[
 \|p\|_{\Gamma,*}^2=
 \int_{\mathbb R}|p(c_*+iy)|^2
 \frac{(2\pi)^{s/2}2^{s/2}}{4\pi\Gamma(s/2)}
 |\Gamma(s/4+iy/2)|^2\,dy,\qquad
 c_*=a_{\mathrm{amp},*}/2-4,
\]
where $s\in\{1,a_{\mathrm{amp},*}\}$ is the fixed original
reference source order. This integral defines a positive definite
norm on $\cP_3$: its density is positive on the real axis and a
nonzero polynomial cannot vanish on that entire axis. Finiteness is
the original WCF2 orthogonal-polynomial norm identity
$h_{n,s}=(2\pi)^{s/2}n!(s/2)_n$, applied to the finite expansion
of the polynomial in that frame, with no change of measure.
With these exact maps and norms, define
\[
 P_{W,*}=\Psi_*P\Psi_*^{-1},\qquad
 p_\delta=\Psi_*u_\rho(\delta),\qquad
 Z_{\varepsilon,\eta,\pm}(\delta)
              =\Psi_*q_{\varepsilon,\eta,\pm}(\delta).
\]
The variable $\delta$ in these formulas parametrizes an auxiliary
target path inside this \emph{fixed} operator and its \emph{fixed}
norm. It does not move the reference arithmetic period. Exact
conjugacy gives all eight fibre identities
\[
 P_{W,*}(Z_{\varepsilon,\eta,\pm}(\delta))=p_\delta.
\]
GF25 and continuity of the fixed linear map give the actual norm
asymptote, for every one of the eight branches,
\begin{equation}
 \sqrt\delta\,\|Z_{\varepsilon,\eta,\pm}(\delta)\|_{\Gamma,*}
       \longrightarrow\frac{\|\Psi_*e_a\|_{\Gamma,*}}{2\gamma}>0.
 \tag{GF26}
\end{equation}
The strict inequality follows because $\Psi_*$ is injective and
the original Gamma norm is positive definite.

The target displacement in its original coordinates is exactly
\begin{equation}
 u_\rho(\delta)-u_\rho(0)
 =\delta^2 d_\gamma-\tfrac12\delta^4e_w,
 \quad d_\gamma=(0,1,-1,\tfrac14-\gamma^2)^{\mathsf T},
 \quad e_w=(0,0,0,1)^{\mathsf T}.\tag{GF27}
\end{equation}
Subtract GF20 at zero to verify every coefficient. Therefore
$\delta^{-2}\|p_\delta-p_0\|_{\Gamma,*}
\to\|\Psi_*d_\gamma\|_{\Gamma,*}>0$.
Combining this with GF26 proves the complete original-norm rate
\begin{equation}
 \boxed{\quad
 \|p_\delta-p_0\|_{\Gamma,*}^{1/4}
 \|Z_{\varepsilon,\eta,\pm}(\delta)\|_{\Gamma,*}
 \longrightarrow
 \frac{\|\Psi_*d_\gamma\|_{\Gamma,*}^{1/4}
       \|\Psi_*e_a\|_{\Gamma,*}}{2\gamma}>0.
 \quad}\tag{GF28}
\end{equation}
The limiting point $p_0$ has no finite preimage under $P_{W,*}$,
because GF12 and the exact conjugacy transfer that assertion as
well. For completeness, the entire nonproper-value locus of this
fixed nonlinear receiving map is precisely
$\Psi_*(\{u_0\mathfrak D(u)=0\})$: invertible linear maps preserve
boundedness, escaping sequences, and convergence. This statement
concerns $P_{W,*}$; the original linear conductor remains the
invertible intertwiner defining $\Psi_*$.

GF26--28 avoid using a singular limit of a moving Vandermonde or a
moving conductor norm. Such a limit would be a separate calculation:
the construction of $\Psi_\delta$ itself assumes distinct original
roots, and its original Vandermonde vanishes at $\delta=0$.
Everything asserted here has instead been proved with the reference
operator, its full original coefficients, and its original norm held
fixed. No assertion about RH follows from introducing this nonlinear
map or this auxiliary target path.

\begin{thebibliography}{9}
\bibitem{ES488} The Clankers, \emph{Erd\H{o}s--Straus project
reader}, canonical 488-page deposited source, version 69,
\texttt{erdos\_straus\_project\_reader.tex}, theorem
\texttt{explicit-four-time-Jacobian-collision}, source lines
19484--19588. Recorded source SHA-256:
\texttt{678b3f70ee7d2a6d00f6c588d7778e390118ab4743a5759612dbba6430984c78}.
The map and four marked source points originate in that local
programme source; the global fibre and nonproper-locus calculations
in this note are new derivations and are not retroactively
attributed to the reader or to a human author of another map.
\bibitem{Bridge} Split-Zero programme,
\emph{The ES--Fable inverse correspondence in the original weighted conductor},
\texttt{FABLE\_TO\_ORIGINAL\_CONDUCTOR.tex}, FC26--35 and FC50--55,
20 September 2026. The four-plane construction uses the original
WCF1--6 finite-shift operator and its original Meixner--Pollaczek
Gamma norms; their human orthogonal-polynomial provenance is
retained there. The complete direct Gram proof also appears in
\texttt{ACTUAL\_FOUR\_PLANE\_METRIC\_AND\_CONJUGACY.tex}, R1--8.
\end{thebibliography}
\end{document}
